\scanpage{82}\mspage{12}
\deleted{\unsure{$\xi$}}
\diagram{\begin{tikzcd}[ampersand replacement=\&,row sep=small,execute at end picture={\draw (current bounding box.west)--(current bounding box.north)--(current bounding box.center)--(current bounding box.north east)--(current bounding box.south east);}]
 \&\deleted{\tilde a}\\ \deleted{\tilde x}\arrow[ur,dash]\arrow[dr]\&\\ \&\deleted{\tilde b}\end{tikzcd}}
Il faut interpréter $\widetilde X_{//(\tilde a,\tilde b)}$ \unsure{soit|disons} donc
\diagram{\begin{tikzcd}[ampersand replacement=\&,row sep=small,column sep=small]
 \&\&\tilde a\arrow[dl,dashed]\arrow[dr,dashed]\&\\
 \widetilde{\unsure{x|\xi}}\arrow[urr,"v"]\arrow[drr,"w"']\&\widetilde{\unsure{\xi|\zeta}}\&\&\widetilde{\unsure{\xi'|b|c}}\\
 \&\&\tilde b\arrow[ul,dashed]\arrow[ur]\&
\end{tikzcd}}
au dessus de
\diagram{\begin{tikzcd}[ampersand replacement=\&,row sep=small,column sep=small]
 \&\&a\arrow[dl,dashed,"\lambda"']\arrow[dr,dashed,"\lambda'"]\&\\
 x\arrow[urr,"v"]\arrow[drr,"w"']\&\xi\&\&\xi'\\
 \&\&b\arrow[ul,"\mu"]\arrow[ur,"\mu'"']\&
\end{tikzcd}}
Donc $\tilde x$ \deleted{\unsure{$\in E$}} $E(x)/R(x)$ peut être défini indifféremment par
\[
 \underbrace{\lambda v,\mu w}_{\substack{\rotatebox{-90}{$\in$}\\\mathrm{Hom}(x,\xi)}},\quad
 \underbrace{\lambda'v,\mu'w}_{\substack{\rotatebox{-90}{$\in$}\\\mathrm{Hom}(x,\xi')}}
\]
de $E(x)$ \unsure{:|;|\textquotesingle} (\unsure{disons} leurs images dans $E(x)/R(x)$) i.e. ces quatre éléments sont congrus mod $R(x)$. On le sait déjà pour $\lambda v\sim\lambda'v$ et $\mu w\sim\mu'w$, en fait $\tilde x=v^*(\tilde a)=w^{\unsure{\prime\prime|*}}(\tilde b)$, donc il faut exprimer l'égalité
\[\deleted{\unsure{\Leftrightarrow}}\qquad v^*(\tilde a)=w^*(\tilde b)\]
\unsure{ou|donc} encore
\[\boxed{\lambda v\sim\mu w}\]
% archive sheet number: 81

